% Problem-section file following ../_template.tex.
\section{Minimal-support frontier for absolutely maximally entangled states}
\label{sec:problem-42}

\paragraph{Problem.}
Determine, for every integer $d\geq 2$, the largest number $\mathcal N(d)$ of
$d$-level parties that admit an absolutely maximally entangled state of
minimal computational-basis support.  To make this precise, let
$m=\lfloor N/2\rfloor$ and write a normalized state as
$|\psi\rangle=\sum_{\boldsymbol{x}\in[d]^N}c_{\boldsymbol{x}}
|\boldsymbol{x}\rangle$, where $[d]=\{0,\ldots,d-1\}$.  The state is
$\operatorname{AME}(N,d)$ when every subsystem $S\subseteq\{1,\ldots,N\}$
with $|S|\leq m$ has reduced state
\begin{equation}
  \operatorname{Tr}_{S^{\mathrm c}}|\psi\rangle\!\langle\psi|
  =\frac{I_{d^{|S|}}}{d^{|S|}}.
  \label{eq:p42-ame-condition}
\end{equation}
Condition~\eqref{eq:p42-ame-condition} implies the computational-basis
support bound
\begin{equation}
  \bigl|\operatorname{supp}_{\mathrm{comp}}(\psi)\bigr|
  :=\bigl|\{\boldsymbol{x}:c_{\boldsymbol{x}}\neq0\}\bigr|
  \geq d^{m}.
  \label{eq:p42-support-bound}
\end{equation}
Minimal support means equality in \eqref{eq:p42-support-bound}, and the
frontier to be determined is
\begin{equation}
  \mathcal N(d):=\max\bigl\{N\geq2:\text{a minimal-support }
  \operatorname{AME}(N,d)\text{ state exists}\bigr\}.
  \label{eq:p42-frontier}
\end{equation}
Equivalently, for any alphabet $\mathcal A$ of size $d$, the existence event
in \eqref{eq:p42-frontier} is characterized by
\begin{equation}
  \begin{split}
  &\text{a minimal-support }\operatorname{AME}(N,d)\text{ state exists}
  \\
  &\quad\Longleftrightarrow\quad
  \exists\,\mathcal C\subseteq\mathcal A^N:\quad
  |\mathcal C|=d^{\lfloor N/2\rfloor},\qquad
  \min_{\substack{\boldsymbol{x},\boldsymbol{y}\in\mathcal C\\
                   \boldsymbol{x}\neq\boldsymbol{y}}}
  d_{\mathrm H}(\boldsymbol{x},\boldsymbol{y})
  =\left\lceil\frac N2\right\rceil+1.
  \end{split}
  \label{eq:p42-mds-equivalence}
\end{equation}
Thus \eqref{eq:p42-mds-equivalence} asks for general, possibly nonlinear,
MDS codes rather than only linear codes
\sourcecite{ref:p42-goyeneche}{GAL+15},
\sourcecite{ref:p42-bernal}{Ber19}.

\subsection*{Status}

Unsolved

\subsection*{Source}

Bernal formulates the minimal-support AME frontier through its equivalent
maximum-distance-separable-code problem and obtains only a conditional general
determination \sourcecite{ref:p42-bernal}{Ber19}.

\subsection*{Progress}

\begin{itemize}
  \item Goyeneche et al. proved the MDS equivalence
  \eqref{eq:p42-mds-equivalence}, while Bernal proved that, for fixed $d$,
  the admissible party numbers form the initial interval
  $2\leq N\leq\mathcal N(d)$
  \sourcecite{ref:p42-goyeneche}{GAL+15},
  \sourcecite{ref:p42-bernal}{Ber19}.  This reduces the frontier problem to a
  maximum-length question for unrestricted $d$-ary MDS codes, which is not
  known in general.

  \item Bernal obtained the exact small-alphabet values and the prime-power
  lower bound
  \begin{equation}
    \begin{gathered}
      \mathcal N(2)=3,\quad \mathcal N(3)=4,\quad \mathcal N(4)=6,
      \quad \mathcal N(5)=6,\quad \mathcal N(6)=3,\quad \mathcal N(7)=8,\\
      \mathcal N(d)\geq d+1
      \quad\text{for every prime power }d\geq3.
    \end{gathered}
    \label{eq:p42-known-values}
  \end{equation}
  The value $\mathcal N(4)=6=d+2$ in \eqref{eq:p42-known-values} comes from
  the even-field exceptional MDS parameters at combinatorial dimension $3$;
  consequently, equality with $d+1$ cannot hold uniformly over prime powers
  \sourcecite{ref:p42-bernal}{Ber19}.

  \item Let $M(k,d)$ denote the maximum length of any, not necessarily
  linear, $d$-ary MDS code of cardinality $d^k$.  Bernal proved the
  conditional implication
  \begin{equation}
    \left.
    \begin{gathered}
      d\geq8\text{ is a prime power},\qquad
      k_0=\left\lfloor\frac{d+2}{2}\right\rfloor,\\
      M(k_0,d)=d+1
    \end{gathered}
    \right\}
    \quad\Longrightarrow\quad
    \mathcal N(d)=d+1.
    \label{eq:p42-conditional-frontier}
  \end{equation}
  The general MDS conjecture supplies the premise in
  \eqref{eq:p42-conditional-frontier}.  Its possible $d+2$ exceptions for
  $d=2^j$ and $k\in\{3,d-1\}$ do not apply here because
  $3<k_0<d-1$ for $d\geq8$; the missing step is the conjecture for general
  nonlinear codes
  \sourcecite{ref:p42-bernal}{Ber19}.
\end{itemize}

\subsection*{References}

\begin{enumerate}[label={},leftmargin=4.8em,labelsep=0.5em]
  \footnotesize
  \item[\textup{[GAL+15]}]\label{ref:p42-goyeneche}
  D. Goyeneche, D. Alsina, J. I. Latorre, A. Riera, and K. \.{Z}yczkowski,
  ``Absolutely Maximally Entangled States, Combinatorial Designs, and
  Multiunitary Matrices,'' \emph{Physical Review A} \textbf{92}, 032316
  (2015).
  \href{https://doi.org/10.1103/PhysRevA.92.032316}%
       {doi:10.1103/PhysRevA.92.032316};
  \href{https://arxiv.org/abs/1506.08857}{arXiv:1506.08857}.

  \item[\textup{[Ber19]}]\label{ref:p42-bernal}
  A. Bernal, ``On the Existence of Absolutely Maximally Entangled States of
  Minimal Support II,'' \emph{Quantum Physics Letters} \textbf{8}, 1--4
  (2019).
  \href{https://doi.org/10.18576/qpl/080101}{doi:10.18576/qpl/080101};
  \href{https://arxiv.org/abs/1807.00218}{arXiv:1807.00218}.
\end{enumerate}

\subsection*{Comment}

The unresolved task is the exact existence frontier, not the classification
of states realizing it.  Support is measured in the fixed computational
basis used in \eqref{eq:p42-support-bound}, and the equivalent coding problem
allows arbitrary alphabets and nonlinear codes.  The linear MDS conjecture
alone therefore does not settle \eqref{eq:p42-frontier}.

\subsection*{Tag}

Absolutely maximally entangled states; Multipartite entanglement; Qudit systems;
Quantum coding theory; Combinatorics

\subsection*{ID}

\texttt{op\_7a17786328198e24}
